%! Introduction to Functions and Change — Unit 1, Lesson 1.0 — Lesson Plan
\documentclass[10pt]{article}
\usepackage{precalculus-article}
\usepackage{precalculus-boxes}
\usepackage{graphicx}
\graphicspath{{images/}}
\newcommand{\TallMath}[1]{$\displaystyle #1\rule[-1.4em]{0pt}{3.2em}$}
\newcommand{\CourseName}{Precalculus}
\newcommand{\MeetingLength}{60 minutes}
\newcommand{\UnitNumberName}{Unit 1: Functions and Change \quad}
\newcommand{\LessonNumberName}{Lesson 1.0: Introduction to Functions and Change}

\begin{document}
\begin{center}
    {\Large\bfseries \CourseName \\
    {\small\bfseries \UnitNumberName \LessonNumberName}}
\end{center}

\begin{tcolorbox}[colback=lilac, colframe=plum, boxrule=0.9pt, arc=2mm,
    left=2mm, right=2mm, top=1.5mm, bottom=1.5mm]
    \textbf{Primary Objective:} Students will identify the input and output
    quantities in a real-world relationship and say which depends on the other,
    explain in everyday language what makes a relationship a \textit{function}
    (each input produces exactly one output), distinguish the allowed
    many-to-one case from the disqualifying one-to-many case, and read a
    pre-drawn graph or a table to say where a quantity is increasing,
    decreasing, or constant.
    \hfill\textit{\small (Unit introduction --- no CED topic)}
\end{tcolorbox}

\boxguard[22]
\begin{skillbox}[Lesson Flow --- Gradual Release (60 minutes)]{lilac}
\small
\renewcommand{\arraystretch}{1.25}
\begin{tabularx}{\linewidth}{@{} l c X @{}}
\textbf{Phase} & \textbf{Min} & \textbf{What students are doing, and where it lives} \\
\midrule
Warm-Up                       & 5  & Silent, individual spiral review --- warm-up sheet. \\
Hook                          & 3  & Whole-class discussion --- hook box atop the guided notes. \\
\textbf{I Do} --- model       & 10 & Watch and annotate while the teacher thinks aloud --- the INPUT \& OUTPUT and FUNCTION rows: definitions, displays, problems 1 and 5. \\
\textbf{We Do} --- guided     & 20 & Fill the notes together, every problem cold-called --- problems 2--4 and 6--8, then the MANY-TO-ONE and READING GRAPHS rows (9--15; the teacher works 9 and 13). \\
\textbf{You Do} --- independent & 14 & Work alone while the teacher circulates --- the You Do row, problems 16--19. \\
Debrief                       & 5  & Whole-class share-out, the headline, and the unit map on the unit cover. \\
Homework Launch               & 3  & Start homework problems 1 and 7 alone; the teacher reads 7 over shoulders. \\
\end{tabularx}

\smallskip
\textbf{If the clock slips:} the whole release lives in the guided notes, so there is
no activity sheet to raid. Cut \textbf{problem 15} first (the Debrief says its answer
aloud), then \textbf{problem 3}. Protect the MANY-TO-ONE row --- problems 11 and 12
are the crux, and homework problems 7, 10, and 11(c) cannot be done without them.
Never cut the Homework Launch to zero: with no exit ticket it is the last look at the
work before it leaves the room.
\end{skillbox}

\renewcommand{\arraystretch}{1.6}
\begin{skillbox}[Priority Ideas \& Skills]{goldbox}
\begin{tabularx}{\linewidth}{@{}X X@{}}
\textbf{Priority Skills} &
\textbf{Key Understandings} \\
\midrule
Identify the input and output quantities in a contextual relationship
and decide which depends on which. &
\begin{itemize}[leftmargin=1.2em,itemsep=2pt,topsep=0pt]
  \item The input (independent) quantity is the one we choose or track;
        the output (dependent) quantity responds to it.
  \item Ordered pairs list input first, output second.
  \item The ``you-control'' test settles most cases in one question.
\end{itemize} \\
Decide informally whether a rule is a function, and justify the decision. &
\begin{itemize}[leftmargin=1.2em,itemsep=2pt,topsep=0pt]
  \item A function sends each input to \textit{exactly one} output ---
        a dependable machine.
  \item \textbf{Target misconception:} a repeated output breaks the rule. It
        does not --- many-to-one is allowed; one-to-many is not. Count arrows
        \textit{leaving each input}, never arrows arriving at an output.
  \item Reversing the roles can turn a function into a non-function.
\end{itemize} \\
Read qualitative behavior from a pre-drawn graph or a table. &
\begin{itemize}[leftmargin=1.2em,itemsep=2pt,topsep=0pt]
  \item Read left to right: rising, falling, and flat stretches say how the
        output changes as the input moves.
  \item A flat stretch means the output holds steady --- not that nothing is
        happening.
  \item Students interpret pre-drawn graphs; no freehand sketching.
\end{itemize} \\
\end{tabularx}
\end{skillbox}

\begin{skillbox}[{Vocabulary, Concepts, \& Theorems}]{greenbox}
\renewcommand{\arraystretch}{1.4}
\begin{tabularx}{\linewidth}{@{} >{\leavevmode\bfseries\color{plum}}p{0.24\linewidth} X @{}}
Input (independent variable)  & The quantity we choose or track; listed first in an ordered pair. \\
Output (dependent variable)   & The quantity that responds to the input; listed second. \\
Function             & A rule that sends each input to exactly one output. (Informal today; notation in 1.1.) \\
Arrow diagram        & Inputs on the left, outputs on the right, an arrow from each input to its output. \\
Many-to-one          & Two or more inputs sharing one output. Allowed --- still a function. \\
One-to-many          & One input with two different outputs. Not allowed --- not a function. \\
Increasing / decreasing / constant & How an output behaves as the input moves forward: rising, falling, or holding steady. \\
\end{tabularx}
\end{skillbox}

\begin{skillbox}[Activate Prior Knowledge \& Spiral Review (5 min)]{lilac}
    Please see attached warm-up (spiral review).
    Skills reviewed: evaluating an algebraic expression at a value,
    solving a two-step linear equation,
    reading a value from a table and describing its trend,
    recognizing a point on a coordinate axis.
\end{skillbox}

\begin{skillbox}[Hook (3 min)]{lilac}
    \textbf{The trustworthy machine.} A vending machine and a claw machine both take
    your money. Press B4 on the vending machine and you get the same bag of chips
    \textit{every single time}; drop the claw and you get\ldots{} who knows.
    \begin{enumerate}[itemsep=2pt]
        \item Which machine can you \textit{trust}? Why?
        \item What would it mean for the vending machine to ``break the rules''?
    \end{enumerate}
    \textit{Discussion prompt:} mathematics reserves a special name for the
    trustworthy machine --- a \textit{function}: same input, same output, every time.
    Collect a few examples of dependable and undependable ``machines'' from life.
    Do not define anything yet; the hook only establishes that
    \textit{dependability} is the property worth naming. The FUNCTION row's arrow
    diagrams draw these two machines.
\end{skillbox}

\begin{skillbox}[Lesson --- I Do / We Do / You Do]{lilac}
\begin{multicols}{2}

\textbf{I DO --- Part 1: Two Quantities, One Dependable Machine} (10 min)

\textit{Notes rows INPUT \& OUTPUT and FUNCTION: read each definition, mark up the
display, work problems 1 and 5. Teacher works; students watch and annotate.
Nothing is cold-called in this phase.}
\begin{itemize}
  \item Name the unit's whole subject out loud: pairs of quantities linked so that
        one \textit{responds} to the other. Read the INPUT \& OUTPUT sentences.
  \item Mark up the phone-plan display: fill the two label boxes yourself ---
        \textit{input} under minutes talked, \textit{output} under the bill ---
        narrating the you-control test as you do.
  \item Work \textbf{problem 1} aloud, the candle: ``I don't choose the height;
        it responds to how long the candle burns.''
  \item Read the FUNCTION sentences, then label the two arrow diagrams:
        \textit{a function} under the vending machine, \textit{not a function}
        under the claw. Point at drop 1's two arrows.
  \item Work \textbf{problem 5}: each student has exactly one ID --- one arrow
        leaves each input.
\end{itemize}

\textbf{WE DO --- Part 2: Naming and Deciding} (6 min)

\textit{Problems 2--4 and 6--8. Class fills the blanks together; cold-call every
one.}
\begin{itemize}
  \item \textbf{Problems 2 and 3} go fast; on 3, make one student read an ordered
        pair aloud as a sentence.
  \item \textbf{Problem 4} is the swap error said out loud --- ask the class who
        chooses the minutes.
  \item \textbf{Problems 6--8}: for each, ask ``how many arrows leave one
        input?'' before anyone says yes or no. Keep it verbal --- notation is
        deliberately absent today.
\end{itemize}

\textbf{WE DO --- Part 3: Many-to-One} (8 min)

\textit{Notes row MANY-TO-ONE --- the crux. Teacher reads the rule, labels the two
diagrams, works problem 9; the class works 10--12, every one cold-called.}
\begin{itemize}
  \item Read the boldface line twice: \textit{count the arrows leaving each input,
        never the arrows arriving at an output.}
  \item Label the diagrams: twins $\to$ birthday is \textit{many-to-one: allowed};
        shoe size 9 $\to$ two students is \textit{one-to-many: broken}.
  \item Work \textbf{problem 9} (the repeated 5 is fine), then let
        \textbf{problem 10} catch the repeated \textit{input}.
  \item \textbf{Problems 11 and 12 are the crux}: Maya's claim is the misconception
        in a student's mouth, and the flip in 12 breaks a function by reversing
        it. Make the class say both halves back before anyone writes.
\end{itemize}

\columnbreak

\textbf{WE DO --- Part 4: A Graph Tells a Story} (6 min)

\textit{Notes row READING GRAPHS. Teacher reads the sentences, labels segment A,
works problem 13; the class labels B and C and works 14--15, every one
cold-called.}
\begin{itemize}
  \item Read the graph left to right as a story --- rising = riding away, flat =
        resting, falling = returning. Ask for the \textit{sentence} before the
        vocabulary word.
  \item Work \textbf{problem 13}: at $t = 5$ the graph \textit{is} the
        information; no formula exists. Insist that segment B is not ``nothing
        happening.''
  \item \textbf{Problem 15} lands ``each time, one distance'' --- the rule, on a
        graph.
\end{itemize}

\textbf{YOU DO --- Part 5: Tables \& Flips} (14 min)

\textit{Notes row TABLES \& FLIPS, problems 16--19 --- the whole independent phase.
Students work alone; circulate and say nothing a neighbor could say instead.}
\begin{itemize}
  \item \textbf{Problem 16} reads the candle table the way the class read the
        graph and matches its flat stretch to segment B: the representation
        changed and the reading did not.
  \item \textbf{Problem 17} is the crux transferred: the repeated 80 is meant to
        feel wrong for a moment.
  \item \textbf{Problems 18 and 19} are the runner. Distance \textit{is} a
        function of time; time is \textit{not} a function of distance ---
        problem 12's flip, with no diagram to lean on. 19 is the finisher's
        question.
\end{itemize}
\end{multicols}
\end{skillbox}

\boxguard
\begin{skillbox}[Active Monitoring]{redbox}
Circulate during the \textbf{You Do} phase --- the fourteen minutes where students
are working problems 16--19 without you --- and check. The crux is
\textbf{problems 11 and 12}; their You Do transfers are \textbf{17} (a repeated
output) and \textbf{19} (the flip).
\begin{itemize}
  \item \textbf{Common error:} swapping input and output (``the bill determines the
        minutes''). Cold-call: ``Which quantity do \textit{you} control?''
  \item \textbf{Common error:} rejecting many-to-one rules --- writing ``no'' on
        problem 17 because 80 appears twice. Ask: ``Point at the input 2. How many
        scores does it have? Now the input 3.''
  \item \textbf{Common error:} stating the rule about outputs (``each output has
        one input''). Insist on the direction: the rule constrains each
        \textit{input}.
  \item \textbf{Graph and table reading:} describing a flat stretch as ``nothing
        happening,'' or naming single hours instead of an interval in 16. Push
        for ``from hour \underline{\ \ } to hour \underline{\ \ }.''
\end{itemize}
\end{skillbox}

\begin{skillbox}[Differentiation --- During You Do (14 min)]{redbox}
There is no separate activity sheet: students work the notes' You Do row (problems
16--19) alone while you circulate. Differentiate by \textit{where you stand}, not
by handing out different paper.

\begin{itemize}
  \item \textbf{Support.} Sit with a struggling student on \textbf{problem 16} and
        have them put a finger on the height row and say ``down, down, same,
        down'' before writing any interval. If they cannot read the table's
        story, 17 will not land either.
  \item \textbf{On level.} Let the class run at \textbf{problem 17} unaided. Do not
        rescue the repeated 80; if a student is stuck, point at the arrow
        diagrams in the MANY-TO-ONE row and walk away.
  \item \textbf{Extend.} \textbf{Problems 18 and 19} are the finisher's questions
        and need no setup. A student who resolves both halves --- distance is a
        function of time, time is not a function of distance --- presents 19 in
        the Debrief. Hand anyone done early the AP Practice's free response,
        part (d): the same flip on a food truck.
\end{itemize}
\end{skillbox}

\boxguard[20]
\begin{skillbox}[Debrief (5 min)]{redbox}
Whole class, teacher-led, packets still open. Three moves, in order:
\begin{enumerate}[itemsep=3pt]
  \item \textbf{Share out the You Do} (2 min) --- one answer per problem, not a
        full review. Problem 16: the decreasing and constant intervals. Problem
        17: why the repeated 80 leaves score a function of hours. Problems 18 and
        19: the runner, both directions. If a student settled 19, they present
        it --- it is the only exposure the rest of the class gets to it before
        homework problem 10 asks for it alone.
  \item \textbf{Name the headline} (1 min) --- say it, then make the class say it
        back: \textit{a function sends each input to exactly one output; two
        inputs may share an output, but one input may never have two.}
        Homework problem 12 asks for it in writing.
  \item \textbf{Point forward} (2 min) --- hold up the unit cover and read its
        lessons table aloud, without adding to it. Students should leave knowing
        that 1.1 gives functions their notation and that the rest of the unit is
        about how quantities change together.
\end{enumerate}
\textbf{Do not} re-teach here, take new questions, or let the homework start early.
A debrief that runs long eats the Homework Launch.
\end{skillbox}

\begin{skillbox}[Homework Launch (3 min)]{redbox}
\textbf{There is no exit ticket.} The formative read comes twice: circulating the
You Do, and again here. Write the due date on the board, have students copy it into
the cover's Homework row, then start \textbf{homework problems 1 and 7} alone and in
silence (3 minutes). Circulate straight to the students you flagged on problem 17.
\begin{itemize}[itemsep=2pt]
  \item \textbf{Problem 1} is the warm-up to the page: input and output for a phone
        charging. A swap here means the you-control test did not stick.
  \item \textbf{Problem 7} is the diagnostic: Jordan says favorite color is not a
        function because many students picked blue. A student who agrees with
        Jordan still believes a repeated output breaks the rule.
\end{itemize}
\textbf{Sort what you see into three piles} as you walk: \textit{rule solid}
(nothing to do), \textit{agrees with Jordan} (open the 1.1 warm-up window with one
arrow diagram and ``count the arrows leaving each input''), and \textit{input and
output swapped} (a one-minute check-in at the door of 1.1).
\end{skillbox}

\begin{skillbox}[Reinforcement \& Extension]{goldbox}
\textbf{Homework --- printed, graded (2 pages).} Last in the packet; due on the date
the teacher sets. Every context is new. \textit{Part A} (1--3) names input and
output --- a phone charging, a café's hot chocolate sales, and one ordered pair read
as a sentence. \textit{Part B} (4--7) decides function or not: a four-rule table
including two flips, an arrow diagram with two many-to-one outputs, a list of pairs
with a repeated input, and the diagnostic Jordan claim. \textit{Part C} (8--10)
reads a bathtub's depth graph --- all three behaviors, a flat stretch in context,
and the two-direction question. \textit{Part D} (11--12) reads a parking garage
table, flips it in 11(c), and asks for the headline in writing.

\textbf{AP Practice --- extra credit (2 pages).} Optional, before the homework in the
packet. Section I is five AP-style multiple-choice items (four options) whose
distractors are the lesson's errors: city $\to$ ZIP code (1, the one-to-many rule
hiding among functions), a constant fee called ``not a function'' (2), arrows counted
at the output instead of the input (3), reading a graph's intervals (4), and input
and output swapped (5). Section II is one free-response question on a food truck's
running taco total: input and output with units (a), a constant interval in context
(b), many-to-one (c), the flip (d), and why a running total can never decrease (e)
--- the finisher's question.

\textbf{Preview:} Next lesson (1.1, Function Fundamentals) gives functions their
notation --- $f(x)$ --- along with evaluating, domain, and range. The homework's
closing box shows one of today's sentences rewritten as $C(3) = 11$.
\end{skillbox}

\begin{teachernote}[Warm-Up]
Five minutes, silent start, no calculators. The four items are pure Algebra 1
maintenance --- this is also your first diagnostic look at the class. A student
stuck on item 1 or 2 needs watching when evaluation returns with $f(x)$ notation
in 1.1; item 3's trend question previews the READING GRAPHS row.
\end{teachernote}

\begin{teachernote}[Guided Notes]
The notes are one Main Ideas / Notes table, four instruction rows and a You Do row;
in each instruction row you read the printed definition, mark up the display, and
work the first problem of the grid (1, 5, 9, 13), and the class works the rest with
the pen in their hand, every problem cold-called. Keep the first two rows verbal and
fast --- notation is deliberately absent and arrives in 1.1; do not front-load it.
The MANY-TO-ONE row is the one worth slowing for: label both diagrams yourself,
read the boldface arrow-counting sentence twice, and treat Maya's claim (11) and
the flip (12) as the crux, not as practice. The bike-ride graph is pre-drawn;
students label and read, never sketch. The old read-only unit map is no longer in
the notes --- it is read from the unit cover in the Debrief.

The You Do row (16--19) is the entire independent phase and carries fourteen
minutes, so treat it as the lesson's third act. Problem 17 is the crux transferred
to a table; 19 is the deepest idea of the day --- the runner's distance is a
function of time, but time is \textit{not} a function of distance --- and it needs
no algebra, only patience. If the clock slips, cut problem 15 in the We Do rather
than touching the You Do.
\end{teachernote}

\begin{teachernote}[Debrief]
Guard these five minutes --- they are the first thing the clock takes and the last
thing you should surrender. Borrow from problem 15 instead, and never from the
three-minute Homework Launch: with no exit ticket it is the last look you get. The
share-out is one answer per You Do problem so that a student who never reached 19
still hears the runner flip, which homework problems 10 and 11(c) ask for alone. Do
not let the unit map turn into a preview lecture: two minutes of reading aloud,
whose only job is to say that today's two words are the spine of everything that
follows.
\end{teachernote}

\begin{teachernote}[AP Practice]
Extra credit, optional, and not a substitute for the homework --- say so when you
hand out the packet. Score it however extra credit works in your gradebook; the
cover gives it a $+$ cell outside the total. Multiple-choice answers are C, B, C, C,
A. Item 2 is the most useful one to glance at: a student who circles (A) thinks a
constant output is not a function --- the lesson's misconception in its most
extreme form. In the free response, accept any plausible story in (b) and any
correct reasoning in (e) that rests on ``a sale cannot be un-sold''; (d) is problem
19's flip and should be graded for naming the input 90 with two outputs.
\end{teachernote}

\begin{teachernote}[Homework]
Graded, two pages, due on the date you set; students start problems 1 and 7 in the
Homework Launch. If time is short, grade problems 7, 10, and 12 for accuracy and
the rest for completion: 7 (Jordan) and 10 (both directions on the bathtub) are the
two items that say whether the many-to-one rule landed, and 12 is the headline in
the student's own words --- together they predict 1.1 better than Part A does. On
problem 4, ``each teacher $\to$ a class they teach'' draws the most wrong
\textit{yes} answers; a student who says yes is picturing one class per teacher,
which is a fair conversation, not a deduction. On 11(a)--(b), accept intervals
written either way (``1 to 3,'' ``hours 1--3'').
\end{teachernote}
\end{document}
